\documentclass[10pt,a4paper]{amsart}
\usepackage[margin=19mm]{geometry}
\usepackage{amsmath,amssymb,mathtools}
\usepackage[scr=rsfs,cal=euler]{mathalfa}
\usepackage{xfrac}
\usepackage[unicode,hidelinks,bookmarks=false]{hyperref}
\newcommand{\ie}{i.e.\ }
\newcommand{\cf}{cf.\ }
\newcommand{\resp}{respectively\ }
\newcommand{\Subsection}{\subsection}
\providecommand{\nobreakdash}{\nobreak\hbox{-}}
\setlength{\emergencystretch}{2em}
\multlinegap0pt
\usepackage{bm}
\usepackage{bbm}
\usepackage{dsfont}
\usepackage{xspace}
\usepackage{cancel}
\usepackage{mathtools}
\usepackage{amsthm}
\makeatletter
\@ifundefined{lemma}{\newtheorem{lemma}{Lemma}}{}
\makeatother
\newcommand{\C}{\mathbb{C}}
\newcommand{\dprod}{\displaystyle\prod}
\newcommand{\ii}{\mathbf{i}}
\title[Determinants of modular Collatz graphs and variants]{Determinants of modular Collatz graphs and variants}
\author{Achilleas Karras, Benne de Weger}
\date{Source: arXiv:2601.15463v1 (CC BY 4.0)}
\begin{document}
\maketitle

\noindent\emph{Source and licence.} Determinants of modular Collatz graphs and variants. Version arXiv:2601.15463v1 (CC BY 4.0), distributed under the Creative Commons Attribution 4.0 International licence, as recorded in the arXiv licence metadata for this exact version and shown on the abstract page https://arxiv.org/abs/2601.15463v1. Full rights notice: licenses/arxiv-2601.15463-CC-BY-4.0.txt.\par\smallskip

\noindent\textit{Editorial scope.} Counted unit: the Lemma whose source label is \texttt{lem:eigenvaluesX} in the article (source0.tex lines 1322--1325, source label \texttt{lem:eigenvaluesX}), delivered together with the article's own complete original proof of it (source0.tex lines 1345--1394, one \texttt{proof} environment of 50 lines). No mathematics is written, repaired or re-derived: statement and proof are the article's own text line for line. The proof is detached from the statement in the source: the article states the result at lines 1322--1325 and prints its proof at lines 1345--1394, and the proof environment's own header names the statement, so the association is the article's own. Required context retained uncounted: none. Every definition, display and label of those lines is retained unchanged. Omitted from this package and not redistributed: 1--1321, 1326--1344, 1395--1661. Nothing inside a retained excerpt is omitted or altered: every retained body is delivered exactly as the article's lines are written, so no omission receipt of any kind accompanies this package. Citations: the retained excerpts carry no citation command at all, so no bibliography entry of the article is transcribed and no reference is redistributed. Reference adaptation: none was needed and none was made; every reference inside the delivered unit resolves inside it, so no unresolved reference or citation is produced. Printed slips kept literal: no printed word, symbol, hypothesis or number of the counted statement or of its proof was altered. Layout and class: the article's own class is \texttt{article} with packages including amssymb, latexsym, amsmath, epsfig, amsthm, xcolor; the wrapper is the lane's pinned \texttt{amsart} editorial wrapper at 10pt on A4, which restates the theorem-like environments the retained text needs and adds the article's own macro definitions that the retained text needs (\texttt{\textbackslash C}, \texttt{\textbackslash dprod}, \texttt{\textbackslash ii}), with extra wrapper packages (bm, bbm, dsfont, xspace, cancel, mathtools). The delivered numbering is the unit's own. Assets: no retained span contains a figure environment, a graphics-inclusion command, a PDF-page inclusion, a table or a TikZ picture; no image file is copied into this package, the package declares no asset dependency and no image file is redistributed with it. Licence: version arXiv:2601.15463v1 of this article is distributed under the Creative Commons Attribution 4.0 International licence, as recorded in the arXiv licence metadata for this exact version and shown on the abstract page https://arxiv.org/abs/2601.15463v1. A full rights notice is delivered at licenses/arxiv-2601.15463-CC-BY-4.0.txt. Characters: every retained excerpt is pure ASCII, so the delivered engine, XeLaTeX with its default Latin Modern text font, reports no missing character.
\begin{lemma} \label{lem:eigenvaluesX}
    The characteristic polynomial of $ X_+ + X_- $ is
    $ (x-2) \displaystyle\prod_{d \mid N, d > 1} \left(x^{k_d}-1\right)^{\frac{\phi(d)}{k_d}} $.
\end{lemma}

\begin{proof}[Proof of Lemma \ref{lem:eigenvaluesX}]
    The main idea of the proof is to change to a Discrete Fourier basis. Let $ \zeta = e^{2 \pi \ii / N} $,
    and $ F = \left( \zeta^{ij} \right) _{i,j} $ be the base change matrix from the standard basis of $ \C^N $
    to the Discrete Fourier basis. So we are interested in
    \[ \hat{X} = F^{-1} \cdot (X_+ + X_-) \cdot F \]
    as it has the same spectrum of $ X_+ + X_- $, and it's easier to compute it this way. We have
    \begin{eqnarray*} \left((X_+ + X_-) \cdot F \right)_{i,j}
        & = & (X_+ + X_-)_i \cdot F^{\top}_j = F_{4^{-1}(2i+1),j} + F_{4^{-1}(2i-1),j} \\
        & = & \zeta^{4^{-1}(2i+1)j} + \zeta^{4^{-1}(2i-1)j} = \zeta^{2^{-1}ij} c_j
    \end{eqnarray*}
    where $ c_j = \zeta^{4^{-1}j} + \zeta^{-4^{-1}j} $. We write $ \Pi_2 $ for the permutation matrix of the
    permutation $ \pi_2: n \to 2 n \pmod{N} $, and $ \Delta_c $ for the diagonal matrix with
    $ c_0, c_1, \ldots, c_{N-1} $ on the main diagonal, and we now have that
    $ (X_+ + X_-) \cdot F = \Pi_2^{-1} \cdot F \cdot \Delta_c $. To compute $ \hat{X} $ we multiply on the left
    by $ F^{-1}  $, so we are now interested in what conjugation by $ F $ does to $ \Pi_2^{-1} $. This simply
    gives $ \Pi_2 $, since $ \Pi_2^{-1} \cdot F = F \cdot \Pi_2 $. So we arrive at a very nice expression for
    $ \hat{X} $, namely
    \[ \hat{X} = \Pi_2 \cdot \Delta_c . \]
    In other words, this gives the weighted permutation matrix of $ \pi_2 $, with the
    $ c_0, c_1, \ldots, c_{N-1} $ as weights. The cycle decomposition does not depend on the weights, so it
    is known. For each divisor $ d $ of $ N $ there are $ \dfrac{\phi(d)}{k_d} $ cycles of length $ d $. We
    fix a cycle $ O = (j_0, j_1, \ldots, j_{\ell-1}) $ with length $ \ell = k_d $ and
    $ j_{t+1} \equiv 2 j_t \pmod{N} $ for all $ t \pmod{\ell} $. The corresponding weighted permutation
    matrix is
    \[ X_O = \begin{pmatrix}
        0       & c_{j_1} & 0       & \cdots & 0 \\
        0       & 0       & c_{j_2} & \cdots & 0 \\
        \vdots  & \vdots  & \ddots  & \ddots & \vdots \\
        0       & 0       & 0       & \cdots & c_{j_{\ell-1}} \\
        c_{j_0} & 0       & 0       & \cdots & 0 \\
    \end{pmatrix} . \]
    Then we have $ X_O^{\ell} = P_O I_{\ell} $ with $ P_O = \dprod_{t=0}^{\ell-1} c_{j_t} $. So the eigenvalues
    of $ X_O $ satisfy $ \lambda^{\ell} = P_O $, in other words, the characteristic polynomial of $ X_O $ is
    $ x^{\ell} - P_O $.

    The permutation $ \pi_2 $ has one fixpoint, corresponding to $ O = (0) $ and the weight $ c_0 = 2 $, so
    $ P_O = 2 $. This gives a factor $ x - 2 $ in the characteristic polynomial.

    We now prove that for all other cycles $ P_O = 1 $. Let us write $ z_t = \zeta^{4^{-1} j_t} $ for
    $ t = 0, 1, \ldots, \ell-1 $. Then $ P_O = \dprod_{t=0}^{\ell-1} \left( z_t + z_t^{-1} \right) $.
    Using $ j_{t+1} \equiv 2 j_t \pmod{N} $ we find $ z_{t+1} = z_t^2 $, so $ z_t = z_0^{2^t} $. So
    \[ P_O = \dprod_{t=0}^{\ell-1} \left( z_0^{2^t} + z_0^{-2^t} \right) . \]
    For this product we can apply the telescoping identity
    \[ (w - w^{-1}) \dprod_{t=0}^{\ell-1} \left( w^{2^t} + w^{-2^t} \right) = w^{2^{\ell}} - w^{-2^{\ell}} , \]
    and this is nontrivial precisely because $ d > 1 $ implies $ z_0 \neq \pm 1 $. Finally note that
    $ z_0^{2^{\ell}} = z_0 $ due to the cycle being closed, so we immediately find
    \[ P_O = \dfrac{z_0^{2^{\ell}} - z_0^{-2^{\ell}}}{z_0 - z_0^{-1}} = 1. \]
    So each cycle $ O $ of length $ \ell = k_d $ for $ d > 1 $ contributes a factor $ x^{k_d} - 1 $ to the
    characteristic polynomial, and there are $ \dfrac{\phi(d)}{k_d} $ of those cycles, so this proves the lemma.
\end{proof}

\end{document}
